\documentclass[10pt]{article}
\usepackage[T1]{fontenc}
\usepackage[utf8]{inputenc}
\usepackage[french]{babel}
\usepackage{amsmath,amssymb,mathrsfs,mathpazo}
\usepackage{geometry}
\geometry{a4paper,margin=2cm}
\usepackage{graphicx,array,tabularx,tabularray,colortbl,multirow,multicol,booktabs,enumitem,pifont,ulem,fancybox,diagbox,cancel}
\usepackage{pstricks,pstricks-add,pst-plot,pst-tree,pst-node,pst-text,pst-func,pst-math,pst-bspline,pst-3dplot,pst-eucl,pst-solides3d}
\let\pscalculate\undefined
\usepackage{pst-bezier}
\usepackage{tikz,pgfplots,tkz-tab}
\usetikzlibrary{arrows,arrows.meta,patterns,calc,positioning,trees,decorations.markings,shapes.geometric,shapes.misc,shapes.arrows,intersections,decorations.pathmorphing,decorations.pathreplacing,backgrounds,shadows,babel}
\pgfplotsset{compat=1.18}
\usepgfplotslibrary{fillbetween}
\usepackage[np]{numprint}
\usepackage{listings,makecell,pas-tableur,twoopt}
\let\barre\undefined
\usepackage[pstricks]{tablvar}
\definecolor{darkgreen}{RGB}{0,100,0}
\providecommand{\red}{\color{red}}
\providecommand{\blue}{\color{blue}}
\providecommand{\green}{\color{green}}
\providecommand{\black}{\color{black}}
\providecommand{\white}{\color{white}}
\providecommand{\gray}{\color{gray}}
\providecommand{\cyan}{\color{cyan}}
\providecommand{\magenta}{\color{magenta}}
\providecommand{\yellow}{\color{yellow}}
\providecommand{\violet}{\color{violet}}
\providecommand{\orange}{\color{orange}}
\providecommand{\euro}{\texteuro}
\providecommand{\EUR}{\texteuro}
\providecommand{\ieme}{\textsuperscript{e}}
\providecommand{\degres}{\ensuremath{^\circ}}
\providecommand{\degree}{\ensuremath{^\circ}}
\providecommand{\vv}[1]{\overrightarrow{#1}}
\providecommand{\Coord}[1]{#1}
\providecommand{\decosix}{}
\providecommand{\starredbullet}{\ensuremath{\bullet}}
\pagestyle{empty}
\setlength{\parindent}{0pt}
\usepackage{adjustbox}\DeclareUnicodeCharacter{2236}{:}\DeclareUnicodeCharacter{2190}{\ensuremath{\leftarrow}}
\providecommand{\R}{\mathbb{R}}
\providecommand{\N}{\mathbb{N}}
\providecommand{\Z}{\mathbb{Z}}
\providecommand{\Q}{\mathbb{Q}}
\providecommand{\D}{\mathbb{D}}
\providecommand{\C}{\mathbb{C}}
\providecommand{\e}{{\mathrm{e}}}
\providecommand{\vect}[1]{\overrightarrow{#1}}
\providecommand{\vectt}[1]{\overrightarrow{\mathrm{#1}}}
\providecommand{\vvt}[1]{\overrightarrow{\mathrm{#1}}}
\providecommand{\np}[2][]{{#2~#1}}
\providecommand{\vv}[1]{{\overrightarrow{#1}}}
\renewcommand{\d}{{\mathrm{d}}}
\providecommand{\degres}{{^\circ}}
\providecommand{\degre}{{^\circ}}
\providecommand{\blue}{\color{blue}}
\providecommand{\red}{\color{red}}
\providecommand{\enskip}{\quad}
\providecommand{\llbracket}{\lbrack\!\lbrack}
\providecommand{\rrbracket}{\rbrack\!\rbrack}
\providecommand{\euro}{{\text{€}}}
\providecommand{\no}{{\text{n°}}}
\providecommand{\interval}[2]{{[#1;#2]}}
\providecommand{\intervFO}[2]{{[#1;#2[}}
\providecommand{\dotfill}{\ldots}

\providecommand{\ds}{\displaystyle}
\providecommand{\textding}[1]{#1}
\providecommand{\Oij}{\(\mathrm{O};\vec{\imath},\vec{\jmath}\)}
\providecommand{\Oijk}{\(\mathrm{O};\vec{\imath},\vec{\jmath},\vec{k}\)}
\providecommand{\Ouv}{\(\mathrm{O};\vec{u},\vec{v}\)}
\providecommand{\intervFF}[2]{\ensuremath{[#1;#2]}}
\renewcommand{\euro}{\ifmmode\text{\texteuro}\else\texteuro\fi}
\DeclareUnicodeCharacter{00B0}{\ifmmode^{\circ}\else\textdegree\fi}
\DeclareUnicodeCharacter{20AC}{\ifmmode\text{\texteuro}\else\texteuro\fi}
% Énoncés : Éducation nationale. Transcription et corrigés : APMEP, Denis Vergès et contributeurs. Découpage et classement : CoopMaths, Sébastien Lozano, Éric Elter et contributeurs. Adaptation de la présentation : Maths Explorer.
% CC BY-SA : https://coopmaths.fr/www/about/
% https://coopmaths.fr/alea/static/e3c/2020/tex/e3c_2020_00_sujet52_1_cor.tex
\begin{document}
\noindent\textit{APMEP / CoopMaths. CC BY-SA. Présentation adaptée par Maths Explorer.}\par\medskip
\subsection*{Question 1}

\(3 \times \dfrac{10^n}{2n+1} = 3 \times \dfrac{10^n}{2n \times 2} = 3 \times \left(\dfrac{10}{2}\right)^n \times \dfrac{1}{2} = \dfrac{3}{2} \times 5^n\).

\((u_n)\) est donc une suite géométrique de premier terme \(\dfrac{3}{2}\) (ou 1,5) et de raison 5.

\subsection*{Question 2}

Avec \(\overrightarrow{AB} \begin{pmatrix} 4 \\ 3 \end{pmatrix}\) et \(\overrightarrow{CM} \begin{pmatrix} x + 1 \\ y - 1 \end{pmatrix}\), on a :
\begin{align*}
&M(x\,;\,y) \in \Delta \\
\iff &\overrightarrow{CM} \cdot \overrightarrow{AB} = 0 \\
\iff &4(x + 1) + 3(y - 1) = 0 \\
\iff &4x + 4 + 3y - 3 = 0 \\
\iff &4x + 3y + 1 = 0.
\end{align*}

\subsection*{Question 3}

\begin{align*}
&2\cos(x + \pi) + 1 = 0 \\
\iff &2\cos(x + \pi) = -1 \\
\iff &\cos(x + \pi) = -\dfrac{1}{2}.
\end{align*}

Or on sait que \(\cos\left(\dfrac{2\pi}{3}\right) = -\dfrac{1}{2}\), on a donc :
\[
\cos(x + \pi) = \cos\left(\dfrac{2\pi}{3}\right),
\]

puis :
\[
(x + \pi) = \dfrac{2\pi}{3} \quad \text{ou} \quad (x + \pi) = -\dfrac{2\pi}{3},
\]

soit :
\[
x = -\dfrac{\pi}{3} \quad \text{ou} \quad x = -\dfrac{5\pi}{3}.
\]

La première n'appartient pas à l'intervalle \(\left[0 \,;\, \dfrac{\pi}{2} \right]\) ; c'est donc la seconde :
\[
-\dfrac{5\pi}{3} + 2\pi = -\dfrac{5\pi}{3} + \dfrac{6\pi}{3} = \dfrac{\pi}{3}.
\]

\subsection*{Question 4}

Le dénominateur étant suprérieur ou égal à 1, la fonction est dérivable et en la dérivant comme un quotient :
\[
f'(x) = \dfrac{\e^x (1 + \e^x) - \e^x \times \e^x}{(1 + \e^x)^2} = \dfrac{\e^x}{(1 + \e^x)^2}.
\]

\subsection*{Question 5}

\( f(x) \) est un trinôme du second degré de coefficient principal \(-4{,}5 < 0\) ; sa représentation graphique est une parabole dont la concavité est tournée vers le bas.

On voit que le maximum de \( f \) est obtenu pour \( x = -2 \) et qu'alors ce maximum est égal à \( 4{,}5 \). Les propositions A, B et D sont donc fausses.

On voit que \( f(-5) = f(1) = 0 \) : donc \(-5\) et \(1\) sont les racines du trinôme. On sait que ce trinôme est négatif sauf entre les racines ; c'est bien ce qui est indiqué dans le tableau C.

\end{document}
